% ===========================================================================
\section{Discussion}
\label{sec:discussion} % ===========================================================================

\subsection{What midtraining installed}

The simplest reading that fits all of our Dispatch results is that midtraining installed a prior on each clause rather than on the Charter. \S\ref{sec:r2} shows that \ac{eft} strengthens the clauses it exercises and leaves the others near the control. \S\ref{sec:one_clause} shows that \ac{eft} which contradicts one clause erodes that clause and leaves the others. \S\ref{sec:data} shows that the prior is carried by worked examples of individual adjudications, not by documents that describe the Charter as a whole. Nowhere does the Charter behave as a single object that stands or falls together, even though the model, asked about it, describes it as one (\S\ref{sec:r1}). Python 4 fits the same reading from the other side: its held-out rules are independent habits that the model can express without the \ac{eft} data making them relevant, and they survive.

One way to say this is Bayesian. Midtraining sets a prior over rulesets and \ac{eft} updates it. Midtraining can raise the weight of ``follow the Charter'' relative to ``maximise coin'' when neither contains the other, which is why agreement-only \ac{eft} resolves to the midtrained motivation. It cannot raise the weight of the seven-clause Charter above the five-clause subset that the \ac{eft} data equally supports, unless the Charter is already a unit in the model's hypothesis space. Our corpus did not make it one. We conjecture that only rules written in terms of features the model already represents are in that space, which would also reconcile our results with the positive single-value generalisation of \citet{li2026modelspec}: a value such as pro-affordability is plausibly an existing feature with wide extension, whereas a seven-clause fictional charter is not.

\subsection{Limitations}

Our experiments span a small number of model families (Gemma~3, Gemma~4, GLM-4.5-Air, and the six substrates of the reproduction) and one seed per cell in most places, so we have not fully elicited the performance of midtraining; the seed spread on our primary metric is about 9\,pp. Our post-training pipeline is deliberately minimal: 10\% of OLMo~3's instruction-tuning tokens, no reasoning training and no \ac{rlhf}. The held-out-clause result rests on a fixed split of two clauses, one of which no cell in our grid executes above 30\%. The post-training-method comparison uses grafted rather than midtrained parents. We are not claiming that \ac{rl} destroys midtrained motivations in general; the Dispatch \ac{rl} result has a plausible setting-specific explanation that we have not ruled out. What we do claim is narrower: the published evidence does not establish that these effects survive ordinary variation in the rest of the pipeline, and when we introduce that variation, they frequently do not.

\subsection{Implications for AI safety}
\label{sec:implications}

Overall we view our results as a moderate downwards update on alignment midtraining; the available public evidence does not justify the apparent degree of confidence placed in it by frontier labs. The specific failure matters more than the size of the update. A midtrained model that endorses a constitution in conversation while acting on only the clauses its post-training happened to demonstrate is a model whose alignment cannot be audited by asking it. The remedy in our setting was trivial, since we could add demonstrations of the missing clauses; the point is that one would not know to.

We think more work should go into (i) settings in which stated and enacted dispositions can be measured separately, (ii) the dependence of midtraining on the consistency of later data, since fractions of a percent of contradicting demonstrations undid it here, and (iii) interaction effects with \ac{rl}.
